\chapter{Monitoring \& Drift Detection}
\label{chap:monitoring}

\paragraph{Riassunto del capitolo.} Un modello in produzione si degrada
silenziosamente se la distribuzione dei dati di servizio si discosta da
quella di addestramento. Questo capitolo presenta il modulo di
\textbf{drift detection} sviluppato come contributo originale del
progetto: combina due test statistici complementari -- Kolmogorov–Smirnov
(forma) e Population Stability Index (massa) -- e li espone tramite CLI
ed endpoint REST. Il modulo si appoggia a un \emph{prediction log} su
JSON-Lines per chiudere il ciclo fra produzione e monitoring.

Questo capitolo descrive il modulo di \textbf{drift detection} sviluppato
come contributo originale del progetto. L'obiettivo è rilevare quando la
distribuzione dei dati di produzione si discosta da quella di
addestramento, condizione che invalida silenziosamente le predizioni di
qualsiasi modello supervisionato.

\section{Cos'è la \emph{drift}}

Un modello in produzione assume \emph{implicitamente} che la
distribuzione dei dati che vede sia la stessa di quella su cui è stato
addestrato (\emph{\acrshort{ml} \emph{training/serving skew}}). Quando
questa ipotesi cessa di valere, la qualità delle predizioni si degrada.
La tassonomia classica del \emph{dataset shift} di Quiñonero-Candela
\emph{et al.}~\cite{quinonero2008dataset} distingue tre regimi:
\begin{itemize}
  \item \textbf{Data drift (covariate shift):} cambia $P(X)$ ma
        $P(Y \mid X)$ resta uguale.
  \item \textbf{Concept drift:} cambia $P(Y \mid X)$ a parità di $P(X)$.
  \item \textbf{Prior shift:} cambia $P(Y)$ a parità di $P(X \mid Y)$.
\end{itemize}

Il modulo \codefile{src/monitoring/drift.py} si concentra sul \emph{data
drift} numerico (covariate shift), perché è il caso operativo più
realistico per una pipeline GEIS: una nuova campagna di misura può
introdurre celle in condizioni leggermente diverse (sensori ricalibrati,
PI termico aggiornato, batch produttivo nuovo).

\section{I due test scelti}

Implementiamo due test complementari:

\subsection{Kolmogorov--Smirnov (\acrshort{ks})}

È un test statistico non parametrico a due campioni che confronta le
\emph{CDF empiriche} di reference e current:

\[
D = \sup_{x} |F_\text{ref}(x) - F_\text{cur}(x)|.
\]

Il p-value è calcolato analiticamente da \texttt{scipy.stats.ks\_2samp},
basato sul test classico di Massey~\cite{massey1951ks}. Soglia di
default: $\alpha = 0.05$. Vantaggio: sensibile a \emph{qualunque}
cambiamento di forma (anche bimodale, varianza, ecc.). Svantaggio:
sensibile anche a piccole differenze quando $n$ è grande.

\subsection{Population Stability Index (\acrshort{psi})}

Misura quanto si è ``spostata la massa'' fra i bin. Si binnano i dati
secondo i quantili della distribuzione di riferimento, poi:

\[
\mathrm{PSI} = \sum_{i=1}^{B}(p_i - q_i)\,\log\!\frac{p_i}{q_i},
\]

con $p_i$ e $q_i$ frazioni del bin $i$ in reference e current
rispettivamente. Le soglie $0.10$ e $0.25$ sono la convenzione di settore
introdotta nel \emph{credit-scoring}~\cite{siddiqi2017credit} e oggi
diffusa anche per il monitoring ML:
\begin{itemize}
  \item $\mathrm{PSI} < 0.10$: nessun cambiamento significativo;
  \item $0.10 \leq \mathrm{PSI} < 0.25$: drift minore;
  \item $\mathrm{PSI} \geq 0.25$: drift maggiore.
\end{itemize}

\subsection{Perché entrambi}

KS è bravo a vedere \emph{il cambio di forma} ma ha falsi negativi quando
la distribuzione si trasla con la stessa forma. PSI è bravo a vedere
\emph{spostamenti di massa} fra bin ma ha falsi negativi su distribuzioni
multimodali quando una moda compensa l'altra. Combinandoli (drift se
\emph{almeno uno dei due} fire) si riducono i falsi negativi.

\section{API di programmazione}

Le funzioni esposte da \codefile{src/monitoring/drift.py}:

\begin{lstlisting}[language=Python]
from src.monitoring import detect_dataset_drift, ks_drift, psi_drift

# Test puntuale
stat, p, drift = ks_drift(ref_array, cur_array, alpha=0.05)
psi, drift = psi_drift(ref_array, cur_array, threshold=0.25)

# Report dataset-level
report = detect_dataset_drift(reference_df, current_df,
                              features=["Z_real", "Z_imag"])
print(report.share_drifted, report.drift_detected)
print(report.to_dict())
\end{lstlisting}

Ogni feature drifted in almeno uno dei due test rientra nel verdetto.
Una soglia aggregata pari al $25\,\%$ di feature drifted scatta il flag
``drift\_detected = True''.

\section{Il logger di predizioni}

Per costruire la distribuzione \emph{current} a partire dal traffico
reale, c'è bisogno di un canale di telemetria. Il modulo
\codefile{src/monitoring/prediction\_log.py} è un piccolo store
\emph{append-only} su JSON-Lines:

\begin{lstlisting}[language=Python]
from src.monitoring import log_prediction

log_prediction(
    task="task1",
    inputs={"aging": 2, "temperature": 30.0},
    outputs={"best_model": "Gradient Boosting", "R2": 0.99},
    extra={"client": "lab_jenkins"},
)
\end{lstlisting}

Il file \codefile{logs/predictions.jsonl} cresce in modo incrementale (un
record per riga). I record sono leggibili in batch con
\texttt{read\_predictions()} e ingeribili da \codefile{scripts/monitor.py}
come distribuzione \emph{current} (\texttt{--current-format jsonl}).

\section{Esposizione \acrshort{rest}}

Il bridge FastAPI espone tre endpoint dedicati:
\begin{itemize}
  \item \textbf{POST /api/predictions/log} -- aggiunge un record al log;
  \item \textbf{GET /api/predictions/recent?limit=N} -- ultimi $N$
        record;
  \item \textbf{GET /api/monitoring/drift?soc=X\&feature=Y} -- esegue il
        confronto fra una sottosezione del dataset e il resto, demo per
        l'utente.
\end{itemize}

In produzione il frontend (o un client che integri la pipeline) chiama
\texttt{POST /api/predictions/log} dopo ogni \texttt{POST /api/taskN/run},
allegando contesto (utente, ambiente). Un job batch programmato (es.~con
\texttt{schedule}, \texttt{cron} o un sistema di orchestrazione) lancia
\codefile{python -m scripts.monitor --current logs/predictions.jsonl
--current-format jsonl} con cadenza giornaliera o settimanale.

\section{Esempio di drift report}

Confrontando la metà inferiore del dataset (Aging $\leq 2$) con la metà
superiore (Aging $\geq 3$) -- una situazione \emph{deliberatamente
drifted} -- il modulo produce:

\begin{lstlisting}[language=JSON]
{
  "ks_alpha": 0.05,
  "psi_threshold": 0.25,
  "n_features": 6,
  "n_drifted": 5,
  "share_drifted": 0.83,
  "drift_detected": true,
  "features": [
    {
      "feature": "Z_real",
      "ks_statistic": 0.41, "ks_p_value": 0.00,
      "psi": 0.92,
      "drift": true
    },
    {
      "feature": "Frequency",
      "ks_statistic": 0.01, "ks_p_value": 0.99,
      "psi": 0.00,
      "drift": false
    },
    ...
  ]
}
\end{lstlisting}

Il report mostra che $\mathrm{Re}(Z)$ è dichiarato drifted con altissima
significatività (KS $p \approx 0$, PSI $= 0.92 \gg 0.25$), mentre la
frequenza non drifta -- coerente con il fatto che lo sweep è uguale per
ogni Aging.

\section{Limitazioni del modulo attuale}

\begin{itemize}
  \item Nessuna detection di concept drift -- richiederebbe un canale
        di feedback per i veri label, che nel nostro setup non esiste.
  \item Nessun retraining automatico: il loop si chiude oggi con
        un'azione umana (rilancio di
        \texttt{python -m scripts.train\_all}).
  \item Nessuna stima di robustezza statistica con bootstrap (KS e PSI
        sono punto-stima dei loro veri valori).
\end{itemize}

Per scenari più complessi (drift online su stream ad alta frequenza,
detector basati su modelli, supporto a feature non numeriche) la
sostituzione naturale è una libreria specializzata come Alibi
Detect~\cite{breunig2024alibidetect}: gli stessi pattern adottati nel
nostro modulo (KS, PSI, JSON-Lines log) sono compatibili e il
porting è incrementale.
